\documentclass[11pt]{article}
\usepackage[no-math]{fontspec}
\usepackage[french,english]{babel}
\usepackage{amsmath}
\usepackage{amssymb}
\usepackage{amsthm}
\usepackage{mathrsfs}
\usepackage{tikz-cd}
\usepackage{enumitem}
\usetikzlibrary{babel}

\newcommand{\printedpage}[1]{\marginpar{\footnotesize p.~#1}}
\newcommand{\msp}[1]{\marginpar{\footnotesize [#1]}}
% underlined category names (printed underlined, upright)
\newcommand{\uu}[1]{\underline{\mathrm{#1}}}
% circled labels (printed 1^o, 2^o in a circle)
\newcommand{\circled}[1]{\tikz[baseline=(c.base)]\node[draw,circle,inner sep=1pt](c){#1};}
% undirected link (printed as a plain horizontal line)
\newcommand{\lline}{\mathrel{\relbar\joinrel\relbar\joinrel\relbar}}

\begin{document}

\printedpage{19}
% head: A. Grothendieck : Dérivateurs, Ch. VI, Édition des Universités de Montpellier II et Paris VII
% foot: Version 9 novembre 2001

\[
\begin{tikzcd}[row sep=large, column sep=large]
X_s \arrow[r, "u_*^X"] \arrow[d, "i_s"] & X_t \arrow[d, "i_t"] \\
X/s \arrow[r, "(\mathrm{id}_X)/u", "= u_*^{\Phi(X)}"'] \arrow[ur, bend left=30, shorten <=28pt, shorten >=12pt, "\alpha"] & X/t
\end{tikzcd}
\]

\msp{page 31}
So I have to pack up! One can regard $i_X$ at best as an arrow of $\uu{Fib}_0\, S$, and taking the composites one finds an arrow in $\uu{Fib}_0\, S$:
\[
\begin{tikzcd}
X \arrow[r, "i_X"] \arrow[rr, bend right=25, "\beta_X(=\alpha_X)"'] & \Phi(X) \arrow[r, "\Phi(j_X)"] & \Phi C(X),
\end{tikzcd}
\]
i.e. (now regarding $\Phi_0$ as a functor with values in $\uu{Fib}_0\, S$) a homomorphism of functors
\[
\beta : \mathrm{id}_{\uu{Fib}_0} S \longrightarrow \Phi_0 p_0, % [sic] \mathrm{id}_{\uu{Fib}_0} S -> \mathrm{id}_{\uu{Fib}_0 S}, the S is printed on the baseline, not in the subscript
\]
where $p_0 : \uu{Fib}_0 S \hookrightarrow \mathrm{Cat}/S$ is the canonical inclusion. This $\beta$ is in $W$ argument by argument, and likewise for $\alpha$ page 29, interpreted this time as
\[
\mathrm{id}_{\mathrm{Cat}/S} \longrightarrow p_0\Phi_0.
\]
What we have proved in the end is this:

\textbf{Packing-up theorem.} \textit{Let $\mathcal{F}$ be a full subcategory of $\mathrm{Cat}/S$ which contains one of the three full subcategories $\uu{Fib}_0 S$, $\uu{Parf}\, S$, $\uu{Cofib}_0 S$, and equipping $\mathcal{F}$ with the ``coarse'' localizer $W$ (induced by that of $\mathrm{Cat}/S$). Then the inclusion}
\[
\mathcal{F} \longrightarrow \mathrm{Cat}/S
\]
\textit{induces an equivalence on the localizers.}

\msp{page 32}
On the other hand, the localized categories of the four categories of split fibrations
\[
\uu{Fibsc}\, S, \quad \uu{Cofibsc}\, S, \quad \uu{Fibsc}_0 S, \quad \uu{Cofibsc}_0 S
\]
escape our arguments. \textit{And for good reason!} It is completely open, over $S$ connected say (e.g. $S = \mathbf{\Delta}^1$), whether for a relative homotopy type $\xi$ over $S$, the induced types $\xi_s$ ($s \in S$) form a local system on $S$, with values in $\mathrm{Hot}_W$, as would be the case if the theorem announced at the beginning were true. On the contrary, \textit{every} arrow of $\mathrm{Hot}_W$ must come from an object of $\mathrm{Hot}_W(\mathbf{\Delta}^1)$. The localizations of the 9 other categories of the diagram p. 8, and of all those discussed in the theorem we have just proved, must correspond ``visibly'' to the full subcategories of $\mathrm{HOT}_W(S)$ formed by the locally constant HOT-coefficients on $S$.

Also, I realize that it was absurd to expect $\uu{Fibsc}\, S$ and $\uu{Cofibsc}\, S$ to give canonically equivalent localized categories. For the elements $\xi$ of the first
\printedpage{20}
% head: A. Grothendieck : Dérivateurs, Ch. VI, Édition des Universités de Montpellier II et Paris VII
% foot: Version 9 novembre 2001
\msp{page 33}
of the two categories each give rise to a contravariant functor on $S$
\[
\bar{\xi} : S^o \longrightarrow \mathrm{Hot}_W,
\]
those $\eta$ of the second (which is none other than $(\uu{Fibsc}(S^o))^o$) each give rise to a \textit{covariant} functor
\[
\bar{\eta} : S \longrightarrow \mathrm{Hot}_W,
\]
which is not the same! It is only ``the same'' if these are local coefficients, which means exactly that $\xi$ is in $\uu{Fibsc}_0 S$ resp. $\eta$ in $\uu{Cofibsc}_0 S$. So one does indeed expect \circled{$1^o$} that the arrows
\[
\uu{Fibsc}_0 S \longrightarrow \uu{Fib}_0 S, \quad \uu{Cofibsc}_0 S \longrightarrow \uu{Cofib}_0 S
\]
induce equivalences on the localizations (something I was unable to prove a little while ago) (\setcounter{footnote}{15}\footnote{The quasi-inverses being well and truly given by $\Phi_0$, $\Psi_0$ - for [?] by what else then ???}). On the other hand, it is worse than unreasonable to expect the same to hold for the functors
\[
\uu{Fibsc}_0 S \longrightarrow \uu{Fibsc}\, S, \quad \uu{Cofibsc}_0 S \longrightarrow \uu{Cofibsc}\, S,
\]
that should give by localization the fully faithful inclusion functors, of the categories of locally constant objects into $\mathrm{HOT}_W(S)$ resp. into $(\mathrm{HOT}_W(S^o))^o$. And likewise for
\msp{page 34}
the functors induced on the localizations by
\[
(*) \qquad \uu{Fibsc}\, S \longrightarrow \uu{Fib}\, S, \quad \uu{Cofibsc}\, S \longrightarrow \uu{Cofib}\, S.
\]
Since the localizations of the target categories of these two functors look very much as though they identify with the preceding subcategories of locally constant objects, it would therefore seem that: \circled{$2^o$} The preceding functors define retractions of the category $\mathrm{HOT}_W(S)$ resp. $\mathrm{HOT}_W(S^o)^o$ onto the subcategory of locally constant objects on $S$, the localized functors of the inclusion functors
\[
\uu{Fibsc}_0 S \longrightarrow \uu{Fibsc}\, S, \quad \uu{Cofibsc}_0 S \longrightarrow \uu{Cofibsc}\, S
\]
being fully faithful and having as essential images the subcategories in question.

This retraction seems to me a wholly mysterious thing. This is tied to the fact that for the very special derivator $\mathbf{D} = \mathrm{HOT}$, when one has a $\xi \in \mathbf{D}(S)$, one deduces from it (up to $\mathrm{HOT}/S$-equivalence) a type of diagram $X$ over $S$,
\msp{page 35}
which one can then regard (via the construction $C$) as defining a relative homotopy type \textit{locally constant} over $S$.

\printedpage{21}
% head: A. Grothendieck : Dérivateurs, Ch. VI, Édition des Universités de Montpellier II et Paris VII
% foot: Version 9 novembre 2001

So to elucidate the situation, i.e. the relations between the 13 localized categories coming from the diagram on page 8, it remains to see whether the hypothetical assertions I have just made are indeed valid, namely

\begin{enumerate}[leftmargin=*]
\item[\circled{$1^o$}] The functor deduced from
\[
\uu{Fibsc}_0 S \longrightarrow \uu{Fib}_0 S
\]
by passing to the localizations is an equivalence of categories (and in this question one can replace $\uu{Fib}_0 S$ by any other full subcategory of $\mathrm{Cat}/S$ containing it - it amounts to the same thing), and

\item[\circled{$2^o$}] The functor deduced from
\[
\uu{Fibsc}_0 S \longrightarrow \uu{Fibsc}\, S
\]
by localization is fully faithful. (That its image is exactly made up of the locally constant objects of the first localization, interpreted as $\mathrm{HOT}_W(S)$, is tautological.)
\end{enumerate}

There is no longer any need to worry about the side of
\msp{page 36}
cofibered categories, since all one will be able to say about it comes down to what I have just spelled out, for $X^o$ instead of $X$.

In the direction of \circled{$1^o$}, our preceding arguments by dint of $\Phi_0$, $p$ already give us a useful piece of information. But it will be more convenient now to drop $\uu{Fib}_0 S$ or rather, to replace it by $\uu{Fib}\, S$ (which is legitimate, we said so in \circled{$1^o$}), so that we now have a commutative diagram (solid lines)
\[
\begin{tikzcd}
\uu{Fibsc}\, S \arrow[dr, "j"'] & & \uu{Fibsc}_0 S \arrow[ll, hook', "i"'] \arrow[dl, "k"'] \\
& \uu{Fib}\, S, \arrow[ur, dotted, bend right=40, "\Phi_0=\psi"'] &
\end{tikzcd}
\]
where $i$ is a (strictly) fully faithful inclusion, $j$ (hence also $k$) faithful. Finally, we have a canonical functor $\psi$ ($= \Phi_0$) (dotted). The localizers on the top line are the fine localizers (for $\uu{Fibsc}_0 S$ this is the same as coarse $W$-equivalence), on $\uu{Fib}\, S$ % [sic] les localiseur fins -> les localiseurs fins, missing plural s
on the other hand we take coarse $W$-equivalence. (We know that the localization is the same, up to equivalence, as if one restricts to the full subcategory $\uu{Fib}_0 S$, which is quite remarkable. % [sic] "(On sait ... remarquable." -> "(On sait ... remarquable.)", the opening parenthesis is never closed

\msp{page 37}
All four functors $i, j, k, \Phi_0$ are compatible with the localizers, hence pass to the localizations
\printedpage{22}
% head: A. Grothendieck : Dérivateurs, Ch. VI, Édition des Universités de Montpellier II et Paris VII
% foot: Version 9 novembre 2001
\[
\begin{tikzcd}
\mathrm{HOT}_W(S) \arrow[dr, "\bar{j}"'] & & \mathrm{HOT}^{?}_W(S) \arrow[ll, "\bar{\imath}"'] \arrow[dl, "\bar{k}"'] \\
& \mathrm{HOT}^{\mathrm{lc}}_W(S), \arrow[ur, bend right=40, "\bar{\psi}"'] &
\end{tikzcd}
\]
where the superscript ? in the notation means that I do not know how to interpret this category, unless I prove that $\bar{k}$ is an equivalence, and the superscript lc in the bottom category should suggest the intuition that in any case we are dealing with a sort of ``relative homotopy types'' (in a modified sense) \textit{locally constant} on $S$. The problem is to understand this diagram. The natural conjecture is % [sic] dois suggerer -> doit suggérer, verb agreement and missing accent

\begin{enumerate}[leftmargin=*]
\item[$1^o$)] $\bar{k}$ is an equivalence. (Which implies that $\bar{\psi}$ is the quasi-inverse equivalence, that $\bar{k}$, $\bar{\psi}$ essentially allow one to identify $\mathrm{HOT}^{?}_W$ and $\mathrm{HOT}^{\mathrm{lc}}_W$, so that $\bar{\imath}$ and $\bar{\jmath}$ become functors between $\mathrm{HOT}_W$ and $\mathrm{HOT}^{\mathrm{lc}}_W$.)

\item[$2^o$)] $\bar{\imath}$ is fully faithful (NB the essential image is well understood in any case.)
\end{enumerate}

\msp{page 38}
In the direction of $1^o$), I know for now that one has
\[
\bar{k}\bar{\psi} \simeq \mathrm{id}_{\mathrm{HOT}^{\mathrm{lc}}_W(S)},
\]
so $\bar{k}$ is essentially surjective (on objects), and induces \textit{surjections} on the $\mathrm{Hom}(-,-)$, while $\bar{\psi}$ is injective on the isomorphism classes of objects, and on the $\mathrm{Hom}(-,-)$ (i.e. \textit{faithful}).

If the relations between $\mathrm{HOT}^{?}_W$ and $\mathrm{HOT}^{\mathrm{lc}}_W$ turn out to be impossible to make head or tail of (e.g. $\bar{\psi}$ not even fully faithful, i.e. $\bar{k}$ not even faithful), then the category $\mathrm{HOT}^{?}_W$ deserves to be dumped in the wastebasket, and we will keep only
\[
\begin{tikzcd}[row sep=large]
\mathrm{HOT}_W(S) \arrow[d, shift right=1.5, "\bar{j}"'] \\
\mathrm{HOT}^{\mathrm{lc}}_W(S), \arrow[u, shift right=1.5, "\bar{\psi}'"']
\end{tikzcd}
\]
where we define $\bar{\psi}'$ as
\[
\bar{\psi}' = \bar{\imath}\bar{\psi} = \overline{i\psi}.
\]
And the question (more modest) now becomes:

\begin{enumerate}[leftmargin=*]
\item[$3^o$)] Is $\bar{\psi}'$ fully faithful? (The \textit{essential} image is the same as that of $\uu{Fib}_0$ ($\approx$ $\uu{Fib}$), it is made up, unless I am mistaken, of the locally constant objects - to be checked if need be \dots)
\end{enumerate}

\printedpage{23}
% head: A. Grothendieck : Dérivateurs, Ch. VI, Édition des Universités de Montpellier II et Paris VII
% foot: Version 9 novembre 2001

\msp{page 39}
It is hard to imagine that $\bar{k}$, or what amounts to the same thing $\bar{\psi}$, could be an equivalence, without there being at least a ``$W_S$-homotopy''
\[
\mathrm{id}_{\uu{Fibsc}_0 S} \ \leftrightarrow\ \leftrightarrow^{h}\!\leftrightarrow\ \leftrightarrow\ \psi k
\]
in the category of endofunctors of $\uu{Fibsc}_0 S$. To tell the truth, I now realize that the description of $\bar{\psi}$ starting from $\psi : X \longmapsto \Phi_0 X = \Phi C X$ is needlessly complicated, for we had made the detour through $C(X)$ solely to be sure of landing in $\uu{Fibsc}_0 S$, not just in $\uu{Fibsc}\, S$. But if we drop $\uu{Fibsc}_0 S$, we don't give a damn, and the functor $\Phi : X \longmapsto \Phi(X)$ is much more natural. And if we do not drop it, there is no harm in replacing the source $\uu{Fib}\, S$ by $\uu{Fib}_0 S$ (which has the same localization, up to equivalence), and on $\uu{Fib}_0 S$ again there is no need to burden ourselves with $CX$, for the pleasure of applying $\Phi$ to it, since $\Phi(X)$ itself is already in $\uu{Fibsc}_0 S$.

So, the question is this: If $X$ varies
\msp{page 40}
% layout: the arrows nearrow, uparrow, nwarrow are printed under each "Fibsc_0 S", pointing up to it; the arrow labelled c from X to Phi(X) is a long line with its arrowhead printed in the middle
\[
\begin{array}{l}
\text{in } \underset{\nearrow\uparrow\nwarrow}{\uu{Fibsc}_0 S}, \quad \text{is there a homotopy in } \underset{\nearrow\uparrow\nwarrow}{\uu{Fibsc}_0 S}, \quad \text{``natural'', i.e. functorial in } X, \\[2ex]
\qquad\qquad X \xrightarrow{\qquad c \qquad} \Phi(X) = \ \text{Split fibered category with fibers the } s\backslash X
\end{array}
\]
But a path which is made up of arrows
\[
X = X_0 \lline X_1 \lline X_2 \cdots \lline X_n = \Phi(X)
\]
which are all cartesian (and moreover, of course, in $W$!). This therefore means that we would have split fibered categories over $S$, $X_1$, \dots, $X_{n-1}$, depending \textit{functorially} on $X$, which connect
\[
X = X_0 \qquad \text{to} \qquad \Phi(X) = X_n
\]
i.e.
\[
X_s \qquad \text{to} \qquad s\backslash X \qquad \text{???}
\]
I see nothing on the horizon.

\end{document}
